% © 2026 Ing. David Jaroš — CC BY-NC-ND 4.0
% UBT-AI-PROVENANCE-BEGIN
% schema: ubt-ai-provenance/v1
% tier: C_working
% ai_assistance: disclosed
% human_review: risk-based
% editorial_responsibility: Ing. David Jaroš
% policy: ../../AI_PROVENANCE.md
% notice: Working material; exhaustive human review is not claimed.
% UBT-AI-PROVENANCE-END

\documentclass[11pt,a4paper]{article}
\usepackage[margin=2.5cm]{geometry}
\usepackage{amsmath,amssymb,amsthm,mathtools}
\usepackage[most]{tcolorbox}
\newtheorem{theorem}{Theorem}
\newtheorem{corollary}[theorem]{Corollary}
\newtheorem{remark}[theorem]{Remark}
\title{Action-Origin Obstruction for the Canonical UBT Generalized-Dirac Programme}
\author{Ing. David Jaro\v{s}}
\date{20 August 2026}
\begin{document}
\maketitle

\begin{abstract}
The active UBT geometry uses one field only,
$E_\mu=\mathcal N_0^{-1/2}D_\mu\Theta$, with the metric and curved Clifford
matrices derived from the same $E_\mu$.  A previous exact theorem shows that a
first-order generalized-Dirac constraint can preserve metric rank ten without
extra variables when its Jacobian with respect to the original field value is
invertible.  This note addresses the next logical gap: can such a first-order
equation be the Euler--Lagrange equation of the currently documented canonical
quadratic kinetic action?  The answer is no, unless a degeneracy or additional
variational mechanism is derived.  This is an order-of-equations obstruction,
not a rank obstruction.
\end{abstract}

\section{Setup}
Let $\Phi^A$ denote the eight real components of the one biquaternionic field
$\Theta$.  Consider a first-derivative Lagrangian density
\[
L(\Phi,\partial\Phi)=\frac12
G^{MN}H_{AB}(\Phi)\,D_M\Phi^A D_N\Phi^B-V(\Phi),
\]
where $G^{MN}$ is nondegenerate on the chosen spacetime/fibre branch and
$H_{AB}$ is the nondegenerate real field-space pairing induced by the canonical
biquaternionic kinetic form.  Gauge and frame connections may contribute lower
order terms; they do not change the principal derivative-order argument below.

\section{Principal Euler--Lagrange order}
\begin{theorem}[Nondegenerate quadratic kinetic term gives a second-order field equation]
Assume the first-derivative Hessian
\[
W^{MN}{}_{AB}:=\frac{\partial^2L}
{\partial(D_M\Phi^A)\,\partial(D_N\Phi^B)}
=G^{MN}H_{AB}
\]
is nonzero, and in particular nondegenerate in the combined derivative/field
indices on the branch under consideration.  Then the Euler--Lagrange equation
has principal part
\[
-G^{MN}H_{AB}\,D_MD_N\Phi^B,
\]
and is genuinely second order.  It therefore cannot be identical, as a local
differential equation on the same variables, to a genuinely first-order
Dirac equation
\[
\mathcal D(\Phi,D\Phi)=0
\]
with nonvanishing first-order principal symbol.
\end{theorem}

\begin{proof}
For a first-derivative Lagrangian the Euler--Lagrange operator is
\[
\mathcal E_A(L)=\frac{\partial L}{\partial\Phi^A}
-D_M\left(\frac{\partial L}{\partial(D_M\Phi^A)}\right).
\]
Differentiating the momentum term once more, the coefficient of the second
jet $D_MD_N\Phi^B$ is minus the velocity Hessian $W^{MN}{}_{AB}$.  Under the
stated nondegeneracy this coefficient does not vanish, so the differential
order cannot drop to one.  Connection and potential terms alter only lower
order pieces.  A first-order equation has no independent second-jet principal
coefficient, hence the two local differential operators cannot be identical.
\end{proof}

\begin{corollary}[Current canonical scalar kinetic action does not directly yield the generalized-Dirac equation]
The quadratic kinetic action presently documented in the UBT action track,
whose field equation is of $D^\dagger D\Theta=\partial V/\partial\Theta^\dagger$
type, cannot directly have the canonical generalized-Dirac equation as its
Euler--Lagrange equation while the kinetic pairing remains nondegenerate.
\end{corollary}

\section{Why factorisation alone is not enough}
Suppose a second-order operator factorises schematically as
\[
\mathcal P=(\mathcal D-\mathcal M)(\mathcal D+\mathcal M).
\]
Even exact factorisation does not imply
\[
\mathcal P\Theta=0\quad\Longrightarrow\quad
(\mathcal D+\mathcal M)\Theta=0.
\]
The elementary one-dimensional example
\[
(\partial_x-m)(\partial_x+m)f=0
\]
has both $e^{mx}$ and $e^{-mx}$ in its second-order solution space, whereas
$(\partial_x+m)f=0$ selects only the $e^{-mx}$ branch.  Therefore selecting a
Dirac square root requires an independently derived branch-selection
principle, boundary/positivity theorem, or a genuinely first-order degenerate
variational structure.  Merely writing a Clifford factorisation is not such a
derivation.

\section{Canonical Clifford matrices depend on the same first jet}
In the active UBT route,
\[
E_\mu=\mathcal N_0^{-1/2}D_\mu\Theta,
\qquad
\Gamma_\mu=\mathcal C(E_\mu).
\]
Consequently a trial density of the schematic form
\[
L_D=\operatorname{Re}\bigl[\bar\Psi\,
 i\Gamma^\mu(E[\Theta])D_\mu\Psi-\bar\Psi\mathcal M\Psi\bigr]
\]
is not the ordinary Dirac variational problem with externally prescribed
$\Gamma^\mu$.  Because $\Gamma^\mu$ itself depends on $D\Theta$, variation of
$L_D$ contains the chain-rule contribution
\[
\frac{\partial\Gamma^\mu}{\partial(D_\nu\Theta)}\,\delta(D_\nu\Theta).
\]
After integration by parts this generically produces second derivatives.
Thus the usual textbook argument that a Dirac Lagrangian yields a first-order
Euler--Lagrange equation cannot be imported without proving a cancellation or
degeneracy specific to the UBT composite lift.

\section{What would close the gap without extra fields}
There are three logically clean no-extra-field possibilities:
\begin{enumerate}
\item \textbf{Degenerate first-order variational structure.} Derive from UBT a
Lagrangian whose full first-jet Hessian has the precise null structure required
to cancel all second-jet terms while retaining the canonical composite
$\Gamma(E[\Theta])$.
\item \textbf{Second-order action plus derived branch selection.} Keep the
quadratic action, prove an exact factorisation, and derive from the same action,
its boundary conditions, positivity, chirality, or complex-time analyticity a
theorem selecting one Dirac factor.  Factorisation alone is insufficient.
\item \textbf{Equivalent constrained variational principle using only
$\Theta$.} Reformulate the action so that no independent auxiliary field is
introduced and prove equivalence of its Euler--Lagrange solution set with the
desired generalized-Dirac equation.  The equivalence, including the composite
connection and holomorphy constraints, must be proved rather than assumed.
\end{enumerate}

If one of these routes yields a first-order equation
$F(\Psi,e)=0$ with invertible $F_\Psi$, the already proved no-extra-variable
rank theorem applies and preserves metric rank ten.  Hence the immediate open
problem is now sharply separated into two pieces:
\[
\text{(A) variational origin of the first-order equation},
\qquad
\text{(B) rank transversality of that equation}.
\]
Piece (B) has a sufficient exact criterion; piece (A) is not yet closed.

\section{Verification}
\begin{center}
\begin{tabular}{p{0.18\textwidth}p{0.72\textwidth}}
Claim & Verification record \\ \hline
Quadratic kinetic term has nonzero second-order principal coefficient & Analytic theorem above; exact SymPy checker `tools/verify_generalized_dirac_action_order.py`. \\
Factorisation does not imply one first-order factor & Analytic counterexample above; exact SymPy substitution for $e^{\pm mx}$. \\
Lean status & \textbf{LEAN-PENDING}.  The execution environment used for this patch has no `lean` or `lake` executable, and the repository currently exposes no `lean-toolchain`/`lakefile` through code search.  No generated Lean text is claimed as a formal proof.  Formalisation target: the finite-dimensional velocity-Hessian/order lemma and the factorisation counterexample. \\
Scope limitation & The theorem proves an obstruction for the currently documented nondegenerate first-derivative kinetic form.  It does not prove that no special degenerate UBT action can yield the generalized-Dirac equation. \\
\end{tabular}
\end{center}

\begin{tcolorbox}[colback=yellow!7!white,colframe=yellow!55!black,title=Current verdict]
The algebraic rank problem is not the present blocker.  The canonical
nondegenerate quadratic $\Theta$ kinetic action gives a second-order
Euler--Lagrange equation.  A first-order generalized-Dirac equation therefore
requires a derived degeneracy/cancellation, or a rigorously selected factor of
the second-order dynamics.  No new field is required in principle, but the
selection mechanism is a genuine open gap and must not be hidden by treating
$\Gamma(E[\Theta])$ as externally fixed.
\end{tcolorbox}

\end{document}
